% =============================================================
%  Chapter 6 — Interpretive Equilibrium
%  Lean source: FormalAnthropology/IDT_Signal_Ch06.lean
% =============================================================
\chapter{Interpretive Equilibrium}
\label{ch:equilibrium}

\epigraph{%
  ``A convention is a regularity in behavior, sustained by an interest in
  coordination and an expectation that others will do their part.''}%
  {David Lewis, \textit{Convention} (1969)}

% ----------------------------------------------------------------
\section{Introduction: Why Cultures Don't Change}
\label{sec:eq-intro}

Every anthropologist eventually confronts the same puzzle: if cultures are
``arbitrary'' in Saussure's sense---if there is no deep reason for driving
on the left rather than the right, for shaking hands rather than bowing,
for calling a dog \emph{dog} rather than \emph{chien}---then why do cultures
persist at all?  Why doesn't every generation re-randomize its conventions?

The game-theoretic answer, originating with Thomas Schelling's
\emph{Strategy of Conflict} (1960) and formalized by David Lewis in
\emph{Convention} (1969), is deceptively simple: cultural forms persist
because \emph{no individual can profitably deviate}.  They are Nash equilibria
of an implicit coordination game.  A greeting ritual need not be optimal in
any absolute sense; it need only be \emph{locally stable}---better than any
unilateral alternative.

This chapter transplants that insight into the framework of Ideatic Diffusion
Theory.  We define an \emph{interpretation game}---a triple of a repertoire
(the receiver's cultural background), a sender payoff, and a receiver
payoff---and show that the concept of equilibrium has rich structural
consequences when the action space is an ideatic space $\IS$ equipped with
composition $\compose$, void $\void$, depth $\depth$, and resonance
$\res$.

The central results of the chapter fall into three groups.  First,
\emph{existence}: void (silence) is always an equilibrium when it weakly
dominates, and constant-payoff games admit universal equilibria
(\S\ref{sec:eq-babbling}).  Second, \emph{structure}: outcomes preserve
repertoire length, depth is linearly bounded, and void preserves depth
exactly (\S\ref{sec:eq-structure}--\S\ref{sec:eq-depth}).  Third,
\emph{robustness}: equilibria survive payoff agreement, monotone transforms,
and void extension (\S\ref{sec:eq-resonance}--\S\ref{sec:eq-advanced}).

Throughout, we provide anthropological commentary and Lean~4 code listings
that verify every theorem with zero \texttt{sorry} axioms.


% ================================================================
\section{Core Definitions}
\label{sec:eq-defs}

\begin{definition}[Repertoire]
\label{def:repertoire}
A \emph{repertoire} is a finite list of ideas in $\IS$:
\[
  \mathrm{Rep} \;=\; \mathrm{List}\;\IS.
\]
It represents the receiver's stock of cultural concepts, categories, and
schemas---what they bring to any interpretive encounter.%
\footnote{In Bourdieu's terms, the repertoire is the agent's \emph{habitus}
rendered as a finite enumeration of conceptual dispositions.}
\end{definition}

\begin{definition}[Interpretation]
\label{def:interpret}
Given a repertoire $R = [r_1,\dots,r_n]$ and a signal $s\in\IS$, the
\emph{interpretation} of $s$ through $R$ is
\[
  \mathrm{interpret}(R, s) \;=\; [r_1\compose s,\;\dots,\;r_n\compose s].
\]
Each background idea is composed with the incoming signal, producing a list
of new ideas---one per schema the receiver already holds.
\end{definition}

\begin{lstlisting}[caption={Interpretation in Lean~4.}]
def interpret (rep : Repertoire I) (signal : I) : List I :=
  rep.map (fun r => IdeaticSpace.compose r signal)
\end{lstlisting}

\begin{definition}[Interpretation Game]
\label{def:interp-game}
An \emph{interpretation game} $G = (R, u_S, u_R)$ consists of
\begin{enumerate}
  \item a repertoire $R$ (the receiver's cultural background),
  \item a sender payoff $u_S : \mathrm{List}\;\IS \to \mathbb{N}$,
  \item a receiver payoff $u_R : \mathrm{List}\;\IS \to \mathbb{N}$.
\end{enumerate}
Payoffs depend on the \emph{interpretation} (the list of composed ideas),
not on the signal directly.  This captures the anthropological insight that
meaning is always mediated by the receiver's culture.
\end{definition}

\begin{lstlisting}[caption={The interpretation game structure.}]
structure InterpGame (I : Type*) [IdeaticSpace I] where
  repertoire : Repertoire I
  sender_payoff : List I -> ℕ
  receiver_payoff : List I -> ℕ
\end{lstlisting}

\begin{definition}[Outcome]
\label{def:outcome}
The \emph{outcome} of sending signal $s$ in game $G$ is the interpretation
produced by composing the repertoire with $s$:
\[
  \mathrm{outcome}(G, s) \;=\; \mathrm{interpret}(G.\mathrm{Rep},\; s).
\]
\end{definition}

\begin{definition}[Equilibrium]
\label{def:equilibrium}
A signal $s$ is an \emph{equilibrium} of game $G$ if no alternative signal
$s'$ yields a strictly higher sender payoff:
\[
  \mathrm{IsEquilibrium}(G, s) \;\iff\;
  \forall\, s'.\;
  u_S\!\bigl(\mathrm{outcome}(G, s')\bigr)
  \;\le\;
  u_S\!\bigl(\mathrm{outcome}(G, s)\bigr).
\]
This is the Nash equilibrium condition for the sender, given the receiver's
fixed repertoire.%
\footnote{We suppress the receiver's strategic choice in this formalization,
treating the repertoire as exogenous.  Chapter~\ref{ch:multireceiver}
relaxes this.}
\end{definition}

\begin{lstlisting}[caption={Equilibrium in Lean~4.}]
def IsEquilibrium (G : InterpGame I) (s : I) : Prop :=
  forall (s' : I), G.sender_payoff (outcome G s') <=
    G.sender_payoff (outcome G s)
\end{lstlisting}


% ================================================================
\section{Babbling Equilibrium and Constant-Payoff Games}
\label{sec:eq-babbling}

The deepest result of Lewis's \emph{Convention} is often overlooked: in
\emph{every} signalling game, there exists a ``babbling'' equilibrium in
which signals carry no information.  In IDT, this corresponds to the void
signal.

\begin{theorem}[Outcome of Void]
\label{thm:outcome-void}
$\mathrm{outcome}(G, \void) = G.\mathrm{Rep}.$
Silence leaves meaning unchanged: the interpretation of void through any
repertoire is the repertoire itself.
\end{theorem}

\begin{proof}[Proof sketch]
By induction on the repertoire list.  The base case is trivial.  For the
inductive step, $r\compose\void = r$ by the right-identity axiom of
$\IS$, so the head is preserved and the tail follows from the inductive
hypothesis.
\end{proof}

This result is the linchpin of everything that follows.  It says that
silence is \emph{transparent}: the receiver's world is unchanged.  In
anthropological terms, a ritual that communicates nothing leaves the
participants exactly where they were.

\begin{theorem}[Babbling Equilibrium / Schelling's Silence (6.1)]
\label{thm:babbling-equilibrium}
If no signal produces a strictly better payoff than silence, then $\void$
is an equilibrium:
\[
  \bigl(\forall\,s.\; u_S(\mathrm{outcome}(G,s)) \le
  u_S(\mathrm{outcome}(G,\void))\bigr)
  \;\implies\;
  \mathrm{IsEquilibrium}(G,\void).
\]
\end{theorem}

\begin{proof}[Proof sketch]
Immediate from the definition of \textsf{IsEquilibrium}: the hypothesis
\emph{is} the conclusion.
\end{proof}

Lewis (1969) showed that babbling equilibria exist in every signalling game.
In IDT, the proof is one line---because the algebraic structure of $\IS$
(specifically, the existence of $\void$) guarantees a universal fallback.
This is the formal foundation of Schelling's focal-point theory: when
nothing you say can improve on silence, silence is optimal.

\begin{theorem}[Constant-Payoff Equilibrium (6.2)]
\label{thm:constant-payoff-equilibrium}
If every signal yields the same payoff~$k$, then every signal is an equilibrium:
\[
  \bigl(\forall\,s.\; u_S(\mathrm{outcome}(G,s)) = k\bigr)
  \;\implies\;
  \forall\,s.\;\mathrm{IsEquilibrium}(G,s).
\]
\end{theorem}

\begin{proof}[Proof sketch]
For any candidate deviator $s'$, rewrite both payoffs to $k$ and apply
$k\le k$.
\end{proof}

Anthropologically, constant-payoff games formalize \emph{radical cultural
relativism}: if all rituals are equally valued, none faces selection pressure,
and every cultural form persists indefinitely.  The theorem explains why
highly egalitarian societies display persistent ritual variation---there is
no payoff gradient to drive convergence.

\begin{theorem}[Equilibrium is Reflexive (6.3)]
\label{thm:equilibrium-refl}
Every signal is at least as good as itself:
$u_S(\mathrm{outcome}(G,s)) \le u_S(\mathrm{outcome}(G,s))$.
\end{theorem}

This is the trivial direction of the Nash definition---every strategy weakly
dominates itself.  Though mathematically obvious, it serves as the base case
for inductive proofs of equilibrium preservation in later chapters.


% ================================================================
\section{Outcome Structural Properties}
\label{sec:eq-structure}

Outcomes have a rigid structure: signals transform but never create or
destroy ideas.

\begin{theorem}[Outcome Length Invariance (6.4)]
\label{thm:outcome-length}
$|\mathrm{outcome}(G, s)| = |G.\mathrm{Rep}|$ for all signals $s$.
\end{theorem}

\begin{proof}[Proof sketch]
$\mathrm{outcome}$ is a \texttt{List.map}, which preserves length.
\end{proof}

Cultural contact does not change how many schemas a society holds---it
transforms each one.  The ``size'' of a culture's conceptual inventory is
invariant under external signalling.  This is a non-trivial constraint: it
rules out models in which signals can spontaneously generate new categories
in the receiver's mind.%
\footnote{Category \emph{creation} requires mechanisms beyond simple
composition---see the discussion of ``conceptual blending'' in
Chapter~\ref{ch:lossy}.}

\begin{theorem}[Outcome Length Independence (6.5)]
\label{thm:outcome-length-eq}
Two different signals produce interpretation lists of the same length:
$|\mathrm{outcome}(G, s_1)| = |\mathrm{outcome}(G, s_2)|$.
\end{theorem}

An immediate corollary of Theorem~\ref{thm:outcome-length}: both sides
equal $|G.\mathrm{Rep}|$.  Game-theoretically, the ``action space'' of
interpretations has constant dimensionality across all possible signals.

\begin{theorem}[Empty Repertoire Trivial Outcome (6.6)]
\label{thm:empty-repertoire-outcome}
A receiver with no ideas produces no interpretations:
$\mathrm{outcome}(\langle [], \ldots\rangle, s) = []$.
\end{theorem}

A mind with zero cultural content cannot interpret anything.  There is no
``raw perception''---interpretation requires prior ideas.  This is the
formal counterpart of Kant's dictum: ``Intuitions without concepts are
blind.''

\begin{theorem}[Singleton Repertoire Outcome (6.7)]
\label{thm:singleton-outcome}
A single-idea mind produces exactly one interpretation:
$\mathrm{outcome}(\langle [r], \ldots\rangle, s) = [r\compose s]$.
\end{theorem}

Monocultural societies have exactly one interpretation of every event.
The singleton case is the simplest non-trivial example and will serve as
a running test case in Chapters~\ref{ch:persuasion}--\ref{ch:multireceiver}.


% ================================================================
\section{Depth Bounds on Outcomes}
\label{sec:eq-depth}

The depth function $\depth$ measures interpretive complexity.  How does it
behave under the \texttt{outcome} operation?

\begin{theorem}[Outcome Depth Bound (6.8)]
\label{thm:outcome-depth-bound}
\[
  \sum_{i} \depth\bigl(\mathrm{outcome}(G,s)_i\bigr)
  \;\le\;
  \sum_{i} \depth(r_i) \;+\; |R|\cdot\depth(s).
\]
The total depth of all interpretations is bounded by the repertoire's total
depth plus the repertoire size times the signal depth.
\end{theorem}

\begin{proof}[Proof sketch]
Induction on the repertoire.  At each step, the depth of $r\compose s$ is
at most $\depth(r) + \depth(s)$ by the subadditivity of composition depth.
Summing over $|R|$ elements yields the bound.
\end{proof}

Game-theoretically, the sender's ``persuasive reach'' is linearly bounded
by signal complexity times audience size.  This result foreshadows the
persuasion bounds of Chapter~\ref{ch:persuasion}: even before introducing
formal limits on marginal depth, we see that total impact is $O(|R|\cdot\depth(s))$.

\begin{theorem}[Void Outcome Preserves Depth (6.9)]
\label{thm:void-outcome-depth}
$\mathrm{depthSum}(\mathrm{outcome}(G,\void)) = \mathrm{depthSum}(G.\mathrm{Rep}).$
\end{theorem}

Silence preserves total interpretive complexity exactly---not just
approximately, but with zero loss.  The ``cost'' of silence is exactly zero
in every dimension.

\begin{lstlisting}[caption={Outcome depth bound in Lean~4.}]
theorem outcome_depth_bound
    (rep : Repertoire I) (s : I) (sp rp : List I -> ℕ) :
    depthSum (outcome (<rep, sp, rp> : InterpGame I) s) <=
    depthSum rep + rep.length * IdeaticSpace.depth s
\end{lstlisting}


% ================================================================
\section{Resonance and Equilibrium}
\label{sec:eq-resonance}

Resonance ($\res$) captures structural similarity between ideas.  How does
it interact with equilibrium?

\begin{theorem}[Resonant Signals Produce Resonant Outcomes (6.10)]
\label{thm:resonant-signals-outcomes}
If $s_1\res s_2$, then for each background idea $r$ in the repertoire,
$r\compose s_1 \res r\compose s_2$.
\end{theorem}

Structurally similar rituals---Christmas and Hanukkah, weddings across
cultures---produce structurally similar experiences for each participant.
This is why comparative anthropology is possible at all: resonant signals
yield resonant interpretations, ensuring that cross-cultural comparison is
not an act of pure projection.

\begin{theorem}[Self-Resonance of Outcomes (6.11)]
\label{thm:outcome-self-resonance}
Every outcome resonates with itself at each position:
$r\compose s \res r\compose s$.
\end{theorem}

A trivial consequence of the reflexivity of resonance, but conceptually
important: a single signal always produces self-consistent interpretations.
No signal can produce an internally contradictory outcome.

\begin{theorem}[Equilibrium Under Payoff Agreement (6.12)]
\label{thm:equilibrium-payoff-agreement}
If two games $G_1, G_2$ have the same repertoire and their sender payoffs
agree on all outcomes, then an equilibrium for $G_2$ is an equilibrium
for $G_1$:
\[
  \bigl(G_1.\mathrm{Rep} = G_2.\mathrm{Rep}\bigr)
  \;\wedge\;
  \bigl(\forall\,s.\; u_{S,1}(\mathrm{outcome}(G_1,s)) = u_{S,2}(\mathrm{outcome}(G_2,s))\bigr)
  \;\implies\;
  \bigl(\mathrm{IsEquilibrium}(G_2,s) \implies \mathrm{IsEquilibrium}(G_1,s)\bigr).
\]
\end{theorem}

\begin{proof}[Proof sketch]
Substitute the payoff agreement into the equilibrium condition for $G_2$ and
apply it to $G_1$.
\end{proof}

If a cultural form is stable in one payoff environment, it remains stable in
any environment that produces the same relative rankings.  This is the formal
version of the sociological observation that institutions survive ``rebranding''
as long as the underlying incentive structure is preserved.


% ================================================================
\section{Composition and Sequential Signalling}
\label{sec:eq-composition}

Cultural communication is rarely a single signal.  Rituals have multiple
acts; conversations have turns.  This section studies how equilibrium
interacts with sequential composition.

\begin{theorem}[Sequential Signal Composition (6.13)]
\label{thm:sequential-signal}
$(r\compose s_1)\compose s_2 = r\compose(s_1\compose s_2).$
\end{theorem}

A ritual with two acts can be understood as a single compound act---but only
from the perspective of someone who has already absorbed the first act into
their background.  This is the associativity axiom of $\IS$, applied to the
signalling context.

\begin{theorem}[Equilibrium Stability Under Void Extension (6.14)]
\label{thm:equilibrium-void-extension}
$\mathrm{outcome}(G,\; s\compose\void) = \mathrm{outcome}(G,s).$
\end{theorem}

\begin{proof}[Proof sketch]
For each $r$ in the repertoire, $r\compose(s\compose\void) = r\compose s$
by the right-identity axiom.  Hence the map is pointwise equal.
\end{proof}

Padding a message with silence doesn't change the game.  This is the formal
version of ``cheap talk is irrelevant when it carries no information.''  In
anthropological terms, the elaborate silences that frame ritual speech
(the pause before a toast, the breath before a prayer) are formally inert
with respect to the equilibrium structure.

\begin{theorem}[Dual Repertoire Combination (6.15)]
\label{thm:dual-repertoire}
$\mathrm{outcome}(\langle R_1\!+\!+\!R_2, \ldots\rangle, s)
= \mathrm{outcome}(\langle R_1,\ldots\rangle, s)
\!+\!+\;\mathrm{outcome}(\langle R_2,\ldots\rangle, s).$
\end{theorem}

A bilingual mind's interpretation is the union of each language's
interpretation.  There is no mysterious ``third space'' of bilingual
meaning---just the disjoint combination.%
\footnote{This is a strong claim that will be qualified in
Chapter~\ref{ch:echo}, where iterated self-interpretation can create
genuinely new structures through repeated composition.}

\begin{lstlisting}[caption={Dual repertoire combination in Lean~4.}]
theorem dual_repertoire_outcome
    (rep_1 rep_2 : Repertoire I) (s : I) (sp rp : List I -> ℕ) :
    outcome (<rep_1 ++ rep_2, sp, rp> : InterpGame I) s =
    outcome (<rep_1, sp, rp> : InterpGame I) s ++
    outcome (<rep_2, sp, rp> : InterpGame I) s
\end{lstlisting}


% ================================================================
\section{Advanced Equilibrium Properties}
\label{sec:eq-advanced}

We close with three theorems on the robustness of equilibria under
transformations of the payoff structure.

\begin{theorem}[Maximum Payoff Equilibrium (6.16)]
\label{thm:max-payoff-equilibrium}
If a signal $s$ achieves the maximum possible sender payoff~$M$,
then $s$ is an equilibrium:
\[
  u_S(\mathrm{outcome}(G,s)) = M
  \;\wedge\;
  \forall\,s'.\; u_S(\mathrm{outcome}(G,s')) \le M
  \;\implies\;
  \mathrm{IsEquilibrium}(G,s).
\]
\end{theorem}

The most sacred ritual---the one that maximizes collective effervescence
(Durkheim)---is always in equilibrium because nothing can improve on the
maximum.  This is tautological in isolation, but gains force in combination
with depth bounds: since persuasion is bounded (Chapter~\ref{ch:persuasion}),
the maximum payoff is itself bounded.

\begin{theorem}[Equilibrium Preserved Under Monotone Transform (6.17)]
\label{thm:equilibrium-monotone}
If $f:\mathbb{N}\to\mathbb{N}$ is monotone and $s$ is an equilibrium under
payoff~$u_S$, then $s$ is an equilibrium under $f\circ u_S$:
\[
  (a\le b \implies f(a)\le f(b))
  \;\wedge\;
  \mathrm{IsEquilibrium}(G,s)
  \;\implies\;
  \forall\,s'.\; f(u_S(\mathrm{outcome}(G,s'))) \le f(u_S(\mathrm{outcome}(G,s))).
\]
\end{theorem}

Cultural stability is robust under monotone revaluation.  If a society
decides that all payoffs should be doubled (inflation), squared (increasing
returns), or logarithmically compressed (diminishing marginal utility), the
same cultural forms remain stable.  The theorem identifies the essential
feature: it is the \emph{ordering} of payoffs, not their magnitudes, that
determines equilibrium.

\begin{theorem}[Void Equilibrium in Zero-Payoff Games (6.18)]
\label{thm:void-equilibrium-zero}
In a game where every signal yields zero payoff, void is an equilibrium:
$\mathrm{IsEquilibrium}(\langle R, \lambda\_.\,0, \lambda\_.\,0\rangle,\;\void).$
\end{theorem}

\begin{proof}[Proof sketch]
Every signal gives payoff~$0$, so every signal is an equilibrium
(Theorem~\ref{thm:constant-payoff-equilibrium} with $k=0$).  In particular,
$\void$ is.
\end{proof}

In a society where communication carries no reward (Durkheim's \emph{anomie}),
silence is trivially optimal.  This is the formal model of communicative
breakdown: when nothing matters, nothing is said.

\begin{lstlisting}[caption={Void equilibrium in zero-payoff games.}]
theorem void_equilibrium_zero_game (rep : Repertoire I) :
    IsEquilibrium
      (<rep, fun _ => 0, fun _ => 0> : InterpGame I)
      (IdeaticSpace.void : I)
\end{lstlisting}


% ================================================================
\section{Conclusion}
\label{sec:eq-conclusion}

This chapter established the game-theoretic foundations of interpretive
stability within IDT\@.  We showed that the algebraic structure of ideatic
spaces---composition, void, depth, resonance---yields a rich theory of
equilibrium without any additional axioms.

Three themes will recur in the chapters ahead.  First, \emph{void as
universal fallback}: silence is always available and always stable
(Theorem~\ref{thm:babbling-equilibrium}).  Second, \emph{linear depth
bounds}: the impact of any signal is bounded by its own complexity times
the audience size (Theorem~\ref{thm:outcome-depth-bound}).  Third,
\emph{resonance preservation}: structurally similar inputs yield
structurally similar outputs
(Theorem~\ref{thm:resonant-signals-outcomes}).

Chapter~\ref{ch:persuasion} turns from equilibrium existence to the
\emph{limits} of persuasion: how much can a signal change a receiver's
depth, and why propaganda is subadditive.
